\chapter{CONCLUSIONES Y RECOMENDACIONES}
\label{ch:conclusiones}
\thispagestyle{empty}

Este capítulo vincula cada objetivo específico planteado en el Capítulo~\ref{ch:planteamiento} con la evidencia numérica y documental reportada en el Capítulo~\ref{ch:propuesta-desarrollo}. Cada conclusión declara explícitamente su grado de cumplimiento —cumplido, parcial o no alcanzado— y se corresponde directamente con una recomendación operativa y metodológica derivada de sus límites observados. Todas las cifras proceden de los registros auditables de las etapas A a D de la campaña explícita.

\section{CONCLUSIONES}

\textbf{Conclusión 1 (OE1, cumplido).} Se elaboró una especificación física trazable para el arreglo cable--entorno que formaliza el conductor, las capas pasivas del aislamiento y cubiertas, el suelo heterogéneo, las fuentes volumétricas, las condiciones de contorno, las interfaces y las unidades internacionales dentro de contratos registrados con huella SHA-256. Ocho casos del catálogo disponen de referencia FEM refinada en tres niveles de malla; el mayor cambio de $\Tmax$ entre los dos últimos refinamientos fue de \SI{0,078}{\kelvin} y el desequilibrio energético máximo alcanzó el 0,027\,\pct{} (Tabla~\ref{tab:explicit-fem}). La procedencia primaria DC o AC de la resistencia $R_{20}$ se mantiene como requisito previo indispensable para trasladar estos contratos a cables comerciales reales.

\textbf{Conclusión 2 (OE2, parcial).} El artefacto computacional basado en subredes PINN por material se implementó y verificó bajo criterios fijados con anterioridad a la confirmación. La receta congelada con arquitectura \texttt{subdomain} de 32 neuronas y 3 capas por material superó todos los criterios en 7 de las 21 ejecuciones de confirmación realizadas con semillas nuevas (Tabla~\ref{tab:explicit-production}). En detalle, alcanzó aceptación total (3/3) en el cable XLPE homogéneo y en la zona seca lejana, y aceptación parcial (1/3) en la zona seca cercana. Los dos casos plenamente aceptados cumplen tanto la tolerancia explícita del 5\,\pct{} como los umbrales de escenario F3 del plan ($E_T<\SI{0,5}{\kelvin}$ y $\mathrm{NRMSE}<2,5\,\pct{}$), sin alcanzar la cota F2 más exigente ($<\SI{0,1}{\kelvin}$ y $<1\,\pct{}$). No se obtuvo aceptación en el relleno mejorado, los estratos discontinuos, el arreglo plano de Aras ni la arena de Kim, y ninguna comparación mostró insensibilidad al aumento de colocaciones (0 de 8 pares). Por ello, la validez del artefacto queda restringida al dominio delimitado en la Sección~\ref{sec:dominio-validez}.

\textbf{Conclusión 3 (OE3, parcial).} A igualdad de fuente térmica prescrita, las referencias FEM cuantifican el cambio de temperatura máxima generado por cuatro heterogeneidades frente al suelo homogéneo de referencia: \SI{-3,327}{\kelvin} con el relleno mejorado, \SI{+6,009}{\kelvin} con la zona seca cercana, \SI{+0,036}{\kelvin} con la zona seca lejana y \SI{+0,519}{\kelvin} con los estratos horizontales. Aunque la mediana de la PINN concuerda en signo en tres de los cuatro casos (Tabla~\ref{tab:explicit-effects}), ninguna comparación superó la regla de resolución preinscrita (Tabla~\ref{tab:explicit-effect-resolution}). En el relleno y la zona seca cercana falló la aceptación completa de los pares, en la zona seca lejana el efecto FEM (\SI{0,036}{\kelvin}) quedó por debajo de la dispersión entre semillas (\SI{0,131}{\kelvin}), y en los estratos concurrieron ambas limitaciones. En consecuencia, la cuantificación del efecto térmico de la heterogeneidad se sustenta en FEM y no queda validada de forma autónoma mediante el solucionador PINN ni se extiende aún al catálogo paramétrico en trébol ni a la localización continua del punto caliente.

\textbf{Conclusión 4 (OE4, no alcanzado).} La comparación sistemática de ampacidad entre representaciones heterogéneas y homogéneas equivalentes por área mediante el procedimiento acotado de bisección especificado en la Sección~\ref{subsec:busqueda-ampacidad} no se ejecutó en la formulación explícita. Únicamente se completó la verificación del piloto DC dependiente de temperatura en geometría coaxial (PINN \SI{32,575572}{\celsius} frente a FEM \SI{32,571582}{\celsius}; diferencia de \SI{0,003990}{\kelvin}), lo que confirma la ejecución de la ley local sin sustituir la búsqueda de corriente límite. Asimismo, el protocolo preliminar de juicio de expertos definido en el Anexo~\ref{anx:protocolo-expertos} no fue aplicado, y las tablas de ampacidad de campañas reducidas previas no sustentan este objetivo.

\textbf{Contribución y alcance metodológico.} La principal contribución de ciencia del diseño reside en exigir la misma física explícita en PINN y FEM, verificar temperatura y flujo normal en cada interfaz y comprobar el balance con el flujo calculado $-k\nabla T$. Este protocolo separa estrictamente el desarrollo de la confirmación y documenta íntegramente los rechazos observados. Las pruebas efectuadas en régimen FEM 2D para el caso Aras mostraron una diferencia de apenas \SI{0,004}{\kelvin} por redistribución radial debida al efecto piel a igual potencia, pero esta evidencia no autoriza a generalizar los pesos entrenados, validar instalaciones de campo ni afirmar aceleración computacional frente a FEM.

\section{RECOMENDACIONES}

\textbf{Recomendación 1 (derivada de la Conclusión 1).} Documentar para cada cable real analizado su construcción geométrica, el tipo de resistencia declarada (DC o AC), la frecuencia nominal y la temperatura de fuente primaria antes de instanciar el contrato físico. Conservar las huellas SHA-256 de cada especificación y exigir la verificación FEM en tres niveles de malla antes de evaluar cualquier configuración nueva.

\textbf{Recomendación 2 (derivada de la Conclusión 2).} Restringir el uso operativo del artefacto actual a los escenarios homogéneos y de zona seca lejana aceptados en confirmación, aplicando verificación FEM obligatoria en la zona seca cercana y absteniéndose de utilizarlo fuera de ese dominio. Para ampliar la cobertura hacia los casos rechazados, ejecutar la ronda de desarrollo E preinscrita en \path{docs/PROTOCOLO_RONDA_E.md} con semillas completamente nuevas. Dicha iteración debe priorizar la continuidad del flujo en los estratos, el balance en el relleno y la precisión de $\Tmax$ en la zona seca cercana, mientras los casos Aras y Kim requieren rediseñar la arquitectura o el presupuesto.

\textbf{Recomendación 3 (derivada de la Conclusión 3).} Emplear la solución FEM como base primaria para cuantificar los efectos térmicos de la heterogeneidad y reportar estimaciones PINN únicamente cuando superen la regla de resolución preinscrita frente a la dispersión entre semillas. Antes de abordar las veinticuatro combinaciones paramétricas y la geometría en trébol del catálogo extendido, reducir la varianza entre inicializaciones cuando el contraste físico esperado sea comparable al rango entre semillas.

\textbf{Recomendación 4 (derivada de la Conclusión 4).} Ejecutar la campaña de búsqueda acotada de corriente límite implementada en \path{Benchmarks/full_ampacity_study.py} con la receta congelada antes de emitir conclusiones de ampacidad. La evaluación debe cubrir los pares homogéneo, de zona seca cercana y de suelo equivalente por área descritos en la Sección~\ref{subsec:busqueda-ampacidad}, junto con la aplicación del protocolo preliminar de juicio de expertos del Anexo~\ref{anx:protocolo-expertos}.

\textbf{Líneas de trabajo futuro.} Para consolidar extensiones operativas futuras, se recomienda incorporar capacidades caloríficas volumétricas documentadas y pruebas de convergencia temporal en régimen transitorio. También resulta necesario modelar pérdidas dieléctricas, corrientes de pantalla y factores de proximidad en régimen AC multifase. Finalmente, el desempeño comparativo entre PINN y FEM debe evaluarse mediante ejecuciones secuenciales aisladas que contabilicen conjuntamente el tiempo de malla, el entrenamiento y la inferencia.

En suma, la investigación establece un marco verificable y transparente sobre los límites actuales de las redes neuronales informadas por física en la transferencia térmica de cables enterrados, consolidando sus productos trazables en las referencias bibliográficas y los anexos que cierran este documento.
